% Generated from package/output/paper_tables_latex/tab2_descriptives.tex (R pipeline); only layout changed.
\begin{table}[tbp]
\centering
\begin{threeparttable}
\caption{Characteristics of the analysed sample (\nTrials{} trials, \nAth{} athletes).}
\label{tab:desc}
\footnotesize
\setlength{\tabcolsep}{3pt}
\begin{tabular*}{\linewidth}{@{\extracolsep{\fill}}>{\raggedright\arraybackslash}p{0.30\linewidth}rlll@{}}
\toprule
Variable & $n$ & Mean $\pm$ SD & Median (IQR) & Range \\
\midrule
Age at first test (years) & 212 & 16.01 $\pm$ 5.94 & 14.10 (12.00--16.80) & 8.80--37.90 \\
Body mass at first test (kg) & 212 & 53.31 $\pm$ 15.90 & 51.73 (40.86--64.03) & 25.03--95.97 \\
Jump height, flight time (cm) & 805 & 28.19 $\pm$ 8.43 & 26.96 (21.62--34.69) & 8.54--51.47 \\
Concentric peak force (N) & 805 & 1269.26 $\pm$ 420.74 & 1227.01 (923.47--1582.99) & 473.52--3073.22 \\
Concentric peak force ($\times$\,BW) & 805 & 2.35 $\pm$ 0.38 & 2.29 (2.11--2.51) & 1.65--4.70 \\
Eccentric braking impulse (N\,s) & 805 & 27.80 $\pm$ 15.43 & 24.60 (16.59--35.70) & 0.00--95.18 \\
Concentric impulse, first 100 ms (N\,s) & 805 & 62.24 $\pm$ 26.66 & 58.70 (41.76--79.14) & 1.51--187.93 \\
Peak landing force ($\times$\,BW) & 805 & 5.39 $\pm$ 1.97 & 4.97 (3.88--6.67) & 2.19--14.88 \\
Contraction time (s) & 805 & 0.78 $\pm$ 0.19 & 0.77 (0.66--0.90) & 0.31--1.48 \\
Countermovement depth (cm) & 805 & $-$26.55 $\pm$ 9.90 & $-$27.13 ($-$33.55--$-$20.71) & $-$52.39--0.00 \\
Touch-down--lowest COM (s) & 805 & 0.25 $\pm$ 0.11 & 0.24 (0.15--0.31) & 0.07--0.93 \\
Trials per athlete & 212 & 3.80 $\pm$ 2.28 & 3.00 (2.00--5.00) & 1.00--16.00 \\
Trials per session & 327 & 2.46 $\pm$ 0.78 & 2.00 (2.00--3.00) & 1.00--7.00 \\
Sessions per athlete & 212 & 1.54 $\pm$ 0.80 & 1.00 (1.00--2.00) & 1.00--6.00 \\
\bottomrule
\end{tabular*}
\begin{tablenotes}[flushleft]\footnotesize
\item \textit{Note.} Age and body mass are per athlete at the first analysed trial; all other rows are per trial unless stated. Jump height is derived from flight time. Countermovement depth is the vendor estimate of centre-of-mass displacement (negative downwards). Touch-down to lowest COM is the duration of the landing phase in the main analysis. BW = body weight; COM = centre of mass; IQR = interquartile range.
\end{tablenotes}
\end{threeparttable}
\end{table}
